\frontchapter{Danh mục các từ viết tắt}

\begin{longtable}{@{}C{0.9cm}L{3.0cm}L{10.4cm}@{}}
  \toprule
  \textbf{STT} & \textbf{Từ viết tắt} & \textbf{Diễn giải} \\
  \midrule
  \endfirsthead
  \toprule
  \textbf{STT} & \textbf{Từ viết tắt} & \textbf{Diễn giải} \\
  \midrule
  \endhead
  1 & AI & Artificial Intelligence --- trí tuệ nhân tạo \\
  2 & API & Application Programming Interface --- giao diện lập trình ứng dụng \\
  3 & ASR & Automatic Speech Recognition --- nhận dạng tiếng nói tự động \\
  4 & BOM & Byte Order Mark --- dấu nhận diện thứ tự byte ở đầu tệp \\
  5 & CER & Character Error Rate --- tỷ lệ lỗi ký tự \\
  6 & CI & Confidence Interval --- khoảng tin cậy \\
  7 & CRUD & Create, Read, Update, Delete --- tạo, đọc, cập nhật, xóa \\
  8 & CSV & Comma-Separated Values --- dữ liệu phân tách bằng dấu phẩy \\
  9 & DB & Database --- cơ sở dữ liệu \\
  10 & ERD & Entity--Relationship Diagram --- sơ đồ quan hệ thực thể \\
  11 & F1 & Trung bình điều hòa giữa precision và recall \\
  12 & FK & Foreign Key --- khóa ngoại \\
  13 & FR & Functional Requirement --- yêu cầu chức năng \\
  14 & HTML & HyperText Markup Language --- ngôn ngữ đánh dấu siêu văn bản \\
  15 & HTTP & HyperText Transfer Protocol --- giao thức truyền siêu văn bản \\
  16 & IEEE & Institute of Electrical and Electronics Engineers \\
  17 & IQR & Interquartile Range --- khoảng tứ phân vị \\
  18 & JSON & JavaScript Object Notation \\
  19 & KV & Key--Value --- khóa--giá trị \\
  20 & LLM & Large Language Model --- mô hình ngôn ngữ lớn \\
  21 & NFR & Non-Functional Requirement --- yêu cầu phi chức năng \\
  22 & OCR & Optical Character Recognition --- nhận dạng ký tự quang học \\
  23 & OLS & Ordinary Least Squares --- bình phương tối thiểu thông thường \\
  24 & PDF & Portable Document Format --- định dạng tài liệu cố định bố cục \\
  25 & PERFIN & Personal Finance --- tên nguyên mẫu trình bày trong niên luận này \\
  26 & PK & Primary Key --- khóa chính \\
  27 & REST & Representational State Transfer \\
  28 & SQL & Structured Query Language --- ngôn ngữ truy vấn có cấu trúc \\
  29 & STT & Speech-to-Text --- chuyển giọng nói thành văn bản \\
  30 & SUS & System Usability Scale --- thang đo khả dụng hệ thống \\
  31 & TC & Test Case --- mã ca kiểm thử ở mức đặc tả \\
  32 & TTL & Time To Live --- thời gian tồn tại của dữ liệu tạm \\
  33 & UAT & User Acceptance Testing --- kiểm thử chấp nhận người dùng \\
  34 & UI & User Interface --- giao diện người dùng \\
  35 & UML & Unified Modeling Language --- ngôn ngữ mô hình hóa thống nhất \\
  36 & VND & Đồng Việt Nam \\
  37 & WER & Word Error Rate --- tỷ lệ lỗi từ \\
  \bottomrule
\end{longtable}

% Longtable không caption ở trên vẫn tăng bộ đếm; đặt lại để bảng nội dung đầu
% tiên mang đúng số 1.
\setcounter{table}{0}

\frontchapter{Quy ước trạng thái và kết quả}

Báo cáo phân biệt rõ thành phần đã hiện thực, kết quả đã đo và hành vi mới ở
mức thiết kế. Quy ước này ngăn kỳ vọng hoặc smoke test bị trình bày như bằng
chứng chất lượng đã hoàn tất.

\begin{longtable}{@{}L{3.0cm}L{5.2cm}L{5.8cm}@{}}
  \caption*{Quy ước mức độ xác nhận thông tin trong báo cáo}\\
  \toprule
  \textbf{Nhãn} & \textbf{Ý nghĩa} & \textbf{Cách sử dụng} \\
  \midrule
  \endfirsthead
  \toprule
  \textbf{Nhãn} & \textbf{Ý nghĩa} & \textbf{Cách sử dụng} \\
  \midrule
  \endhead
  Đã hiện thực & Có thành phần tương ứng trong mã nguồn, migration hoặc test. & Chứng minh sự tồn tại, chưa tự động chứng minh chất lượng. \\
  Đã đo & Có lệnh, thời điểm và kết quả chạy có thể tái lập. & Được phép đưa số liệu vào phần kết quả. \\
  Mục tiêu & Ngưỡng mong muốn dùng làm tiêu chí nghiệm thu. & Không được trình bày như kết quả đã đạt. \\
  Thiết kế đích & Hành vi dự kiến sau khi hoàn tất sửa lỗi trong phạm vi. & Phải được kiểm thử trước khi chuyển thành ``đã đo''. \\
  Ngoài phạm vi & Không phải mục tiêu của niên luận hiện tại. & Chỉ nêu ở hạn chế hoặc hướng phát triển. \\
  \bottomrule
\end{longtable}
